\documentclass[10pt,a4paper]{amsart}
\usepackage[margin=19mm]{geometry}
\usepackage{amsmath,amssymb,mathtools}
\usepackage[scr=rsfs,cal=euler]{mathalfa}
\usepackage{xfrac}
\usepackage[unicode,hidelinks,bookmarks=false]{hyperref}
\newcommand{\ie}{i.e.\ }
\newcommand{\cf}{cf.\ }
\newcommand{\resp}{respectively\ }
\newcommand{\Subsection}{\subsection}
\providecommand{\nobreakdash}{\nobreak\hbox{-}}
\setlength{\emergencystretch}{2em}
\multlinegap0pt
\usepackage{mathtools}
\usepackage{enumerate}
\usepackage{amssymb}
\usepackage[shortlabels]{enumitem}
\newtheorem{lem}{Lemma}
\newtheorem{prop}{Proposition}
\title{Weak limits of the $J$-flow and the deformed Hermitian-Yang-Mills flow on K{\"a}hler surfaces: boundary cases}
\author{Rei Murakami}
\date{Source: arXiv:2404.00501v2 (CC BY 4.0)}
\begin{document}
\maketitle

\noindent\emph{Source and licence.} Weak limits of the $J$-flow and the deformed Hermitian-Yang-Mills flow on K{\"a}hler surfaces: boundary cases. Version arXiv:2404.00501v2 (CC BY 4.0), distributed by arXiv under the Creative Commons Attribution 4.0 International licence, http://creativecommons.org/licenses/by/4.0/, as recorded in the arXiv licence metadata for this exact version and shown on the abstract page https://arxiv.org/abs/2404.00501v2. Full licence notice: \texttt{licenses/arxiv-2404.00501-CC-BY-4.0.txt}.\par\smallskip

\noindent\textit{Editorial scope.} Counted unit: the article's prop prop (source0.tex lines 453--455). The article's complete original proof of it is delivered with it (source0.tex lines 457--524, one \texttt{proof} environment of 68 lines). No mathematics is written, repaired or re-derived: the statement and the proof are the article's own lines, in the article's own notation, and no retained line is substituted or re-flowed.  Retained uncounted as the source-local context the proof needs: lines 392--394.  Nothing was omitted from any retained span, so the package carries no omission record.  The retained lines cite 4 works, whose entries are transcribed verbatim from the article's own bibliography source in the e-print and delivered with short editorial labels recorded in the item's reference mapping.  No retained line contains a figure environment, an image-inclusion command or drawing code, so no image is delivered and the package declares no asset dependency.  Editorial formatting only, in addition: an AMS \texttt{amsart} preamble that restates the article's own declarations and macros used by the retained lines, namely the environments \texttt{lem}, \texttt{prop} on separate counters, together with the standard packages those lines need.  The retained environments print in this document as Lemma 1, Proposition 1 in order of appearance, whereas the article prints the same items with the numbering of its own full document.  The complete native source of the article as distributed in the arXiv e-print for this exact version is bundled unchanged as \texttt{sources/156972/source0.tex}.
\begin{lem}\label{lemdhymconvexity}
    The $\mathcal{J}$-functional is convex along the dHYM flow with an initial point $\varphi_0$ satisfying $0<\theta_\omega(\alpha_0)<\pi$.
\end{lem}

\begin{prop}
    Let $X$ be a compact K{\"a}hler surface and $\omega$ a K{\"a}hler form. Let $\alpha$ be a smooth closed real $(1,1)$-form. Assume that $0 <\theta_0 < \pi$ and the class $[\alpha-\cot{\theta_0}\omega]$ is nef. If an initial point of the dHYM flow satisfies $0 < \theta_\omega(\alpha_{\varphi_0}) < \pi$, then the time derivative of the dHYM flow converges to zero in $L^2$.
\end{prop}

\begin{proof}
    Note that
    \begin{align}\label{eqdhymderiv}
        \frac{d}{dt}\mathcal{J}(\varphi_t)
        &=\int_X \dot{\varphi_t} \, \mathrm{Im}\left(e^{-\sqrt{-1}\theta_0}(\alpha_\varphi+\sqrt{-1}\omega)^n\right) \notag\\
        &=\sin{\theta_0}\int_X \dot{\varphi}_t \left(\cot{\theta_0}\, \mathrm{Im}(\alpha_\varphi+\sqrt{-1}\omega)^n-\mathrm{Re}(\alpha_\varphi+\sqrt{-1}\omega)^n
        \right) \notag\\
        &=-\sin{\theta_0}\int_X \dot{\varphi}_t^2 \, \mathrm{Im}(\alpha_\varphi+\sqrt{-1}\omega)^n \le0.
    \end{align}
    As in \cite{CL}, let $F_{\omega,-\varepsilon}$ be the operator defined by
    $$F_{\omega,-\varepsilon}(\alpha_\varphi)= \frac{\mathrm{Re}(\alpha_\varphi+\sqrt{-1}\omega)^n-\varepsilon\omega^n}{\mathrm{Im}(\alpha_\varphi+\sqrt{-1}\omega)^n}.$$
    Also, let $\mathcal{J}_{-\varepsilon}$ be the functional defined by
    $$\mathcal{J}_{-\varepsilon}(0)=0, \quad d\mathcal{J}_{-\varepsilon}(\varphi)(\psi)=\int_X \psi \, \left(\cot{\theta_0}-a_0\varepsilon-F_{\omega,-\varepsilon}(\alpha_\varphi)\right)\mathrm{Im}(\alpha_\varphi+\sqrt{-1}\omega)^n,$$
    where $a_0$ is the constant determined by
    $$a_0=\frac{\int_X\omega^n}{\int_X \mathrm{Im}(\alpha+\sqrt{-1}\omega)^n}.$$
    We claim that if the class $[\alpha-\cot{\theta_0}\omega]$ is nef and big, then for any $\varepsilon>0$, the functional $\mathcal{J}_{-\varepsilon}$ is bounded from below on 
    $$\mathcal{H}_\pi=\{ \varphi \in C^\infty(X) \mid  0 < \theta_\omega(\alpha_\varphi) < \pi\}.$$
    Note that the critical point $\varphi$ of $\mathcal{J}_{-\varepsilon}$ satisfies
    $$F_{\omega,-\varepsilon}(\alpha_\varphi)=\cot{\theta_0}-a_0\varepsilon,$$
    which is equivalent to
    \begin{equation}\label{eqMA}
        \left(\alpha_\varphi-(\cot{\theta_0}-a_0\varepsilon)\omega\right)^2=(\left(\cot{\theta_0}-a_0\varepsilon\right)^2+1+\varepsilon)\omega^2.
    \end{equation}
    Since $[\alpha-\cot{\theta_0}\omega]$ is nef and $\omega$ is a K{\"a}hler form, the class $[\alpha-(\cot{\theta_0}-a_0\varepsilon)\omega]$ is K{\"a}hler. By Yau's theorem \cite{Yau}, there exists a unique solution $u_{MA}$ of \eqref{eqMA} such that $\mathrm{Im}\mathrm{CY}_\mathbb{C}(u_{MA})=0$. Here, recall that the functional $\mathrm{CY}_\mathbb{C}$ is given by
    $$\mathrm{CY}_\mathbb{C}(0)=0,\quad d\mathrm{CY}_\mathbb{C}(\varphi)(\psi)=\int_X \psi \, (\alpha_\varphi+\sqrt{-1}\omega)^n.$$
    With the properties of $F_{\omega,-\varepsilon}$ \cite[Lemma 5.6]{GChen}, the same arguments as \cite[Section 4.2]{CL} imply that the twisted dHYM flow
    \begin{equation}\label{eqtdHYMflow}
        \begin{cases}
        \dot{u}_t=F_{\omega,-\varepsilon}(\alpha_u)-\cot{\theta_0}+a_0\varepsilon \\ u(x,0)=u_0(x)\in \mathcal{H}_\pi
        \end{cases}
    \end{equation}
    converges smoothly to $u_{MA}+\mathrm{Im}\mathrm{CY}_\mathbb{C}(u_0)$. Since the flow is a gradient flow of $\mathcal{J}_{-\varepsilon}$, we have
    $$\mathcal{J}_{-\varepsilon}(u_0)\ge\mathcal{J}_{-\varepsilon}(u_{MA})=-C_\varepsilon,$$
    namely, the claim follows.
    Then, along the dHYM flow $\varphi_t$ with the initial point $\varphi_0$ satisfying $0<\theta_\omega(\alpha_{\varphi_0}) < \pi$, we have
    \begin{equation}\label{ineqbdd}
        \mathcal{J}_{-\varepsilon}(\varphi_t)\ge - C_\varepsilon
    \end{equation}
    for some constant $C_\varepsilon$, since 
    $\inf\theta_\omega(\alpha_{\varphi_0})<\theta_\omega(\alpha_{\phi_t})<\sup\theta_\omega(\alpha_{\varphi_0})$
    by \cite[Lemma 3.2]{FYZ}. Recall that the imaginary part of the Calabi-Yau functional $\mathrm{CY}_\mathbb{C}$ 
    is constant along the dHYM flow by \cite[Lemma 2.7]{FYZ}. Thus,
    \begin{align*}
        \mathcal{J}_{-\varepsilon}(\varphi_t)-\mathcal{J}_{-\varepsilon}(\varphi_0)
        &=\int^t_0\int_X \dot{\varphi}_t \, \biggl( (\cot{\theta_0}-\cot{\theta_\omega(\alpha_\varphi)})\mathrm{Im}(\alpha_\varphi+\sqrt{-1}\omega)^n+\varepsilon\, \omega^n\biggr)\\
        &=\frac{1}{\sin{\theta_0}}\mathcal{J}(\varphi_t)+\varepsilon\int_X\varphi_t\, \omega^n\\
        &\le\frac{1}{\sin{\theta_0}}\mathcal{J}(\varphi_t)+\varepsilon C\sup\varphi_t\\
        &\le\frac{1}{\sin{\theta_0}}\mathcal{J}(\varphi_t)+\varepsilon C t,
    \end{align*}
    where we used \cite[Lemma 3.2]{FYZ} again to estimate $\sup\varphi_t$.
    By combining this with \eqref{ineqbdd}, we see that for any $\varepsilon>0$ there exist constants $C$ and $C_\varepsilon$ such that
    \begin{equation}\label{ineqdhymslope}
        \mathcal{J}(\varphi_t)\ge -\varepsilon C t-C_\varepsilon.
    \end{equation}
    By Lemma \ref{lemdhymconvexity}, there exists a constant $\mu\le0$ such that
    $$\lim_{t\rightarrow\infty}\frac{d}{dt}\mathcal{J}(\varphi_t)=\mu.$$
    By dividing both hands of \eqref{ineqdhymslope} by $t$ and letting $t$ tend to $\infty$, we also see that
    $$\lim_{t\rightarrow\infty}\frac{d}{dt}\mathcal{J}(\varphi_t)\ge - \varepsilon C.$$
    Since $\varepsilon>0$ is arbitrary, we get $\mu\ge0$, which concludes $\mu=0$. From \eqref{eqdhymderiv}, we obtain
    $$\int_X \dot{\varphi}_t^2 \, \mathrm{Im}(\alpha_\varphi+\sqrt{-1}\omega)^n \rightarrow0.$$
    By \cite[Lemma 3.2]{FYZ}, we have
    \begin{align*}
        \mathrm{Im}(\alpha_\varphi+\sqrt{-1}\omega)^n
        &=\left(\prod_i\sqrt{\lambda_i^2+1}\right)\left(\sin{\theta_\omega(\alpha_\varphi)}\right)\omega^n\\
        &\ge \inf\left(\sin{\theta_\omega(\alpha_{\varphi_0}})\right) \omega^n.
    \end{align*}    
    Therefore, we conclude the statement.
\end{proof}

\begin{thebibliography}{99}
\bibitem[CL]{CL} J. Chu and M.-C. Lee, \textit{Hypercritical deformed Hermitian-Yang-Mills equation,} preprint, available at arXiv:2107.13192v1.
\bibitem[FYZ]{FYZ} J. Fu, S.-T. Yau, and D. Zhang, \textit{A new flow solving the LYZ equation in Kähler geometry,} J. Differential Geom. \textbf{128} (2024), no. 1, 153--192.
\bibitem[GChen]{GChen} G. Chen, \textit{The J-equation and the supercritical deformed Hermitian-Yang-Mills equation,} Invent. Math. \textbf{225} (2021), no. 2, 529--602.
\bibitem[Yau]{Yau} S.-T. Yau, \textit{On the Ricci curvature of a compact Kähler manifold and the complex Monge-Ampère equation. I.} Comm. Pure Appl. Math. \textbf{31} (1978), no. 3, 339--411.
\end{thebibliography}
\end{document}
